% BEGIN GENERATED ARTIFACT: prf:negation-disjunction-functionally-complete-propositional-logic
\newpage
\phantomsection
\label{prf:negation-disjunction-functionally-complete-propositional-logic}
\LRAProofFor{thm:negation-disjunction-functionally-complete-propositional-logic}

\begin{remark*}[Return]
\hyperref[thm:negation-disjunction-functionally-complete-propositional-logic]{Return to Theorem (Negation and Disjunction are Functionally Complete).}
\end{remark*}

\begin{theorem*}[Negation and Disjunction are Functionally Complete]
The connective basis \(\{\neg,\lor\}\) is functionally complete.
\end{theorem*}

\begin{proof}
\textbf{Professional Standard Proof.}~
TODO: supply the compact proof.
\end{proof}

\begin{proof}
\textbf{Detailed Learning Proof.}~
TODO: supply the detailed learning proof.
\end{proof}

\begin{remark*}[Proof structure]
Use De Morgan's law to define conjunction from negation and disjunction:
\(\varphi\land\psi\equiv\neg(\neg\varphi\lor\neg\psi)\). Then translate any
formula over \(\{\neg,\land,\lor\}\) into a formula over \(\{\neg,\lor\}\).
\end{remark*}

\begin{dependencies}
\begin{itemize}
  \item \hyperref[thm:standard-connectives-functionally-complete-propositional-logic]{Standard Connectives are Functionally Complete}
  \item \hyperref[thm:de-morgan-laws-propositional-logic]{De Morgan Laws for Propositional Logic}
  \item \hyperref[thm:replacement-of-logical-equivalents-propositional-logic]{Replacement of Logical Equivalents}
\end{itemize}
\end{dependencies}

\clearpage
% END GENERATED ARTIFACT: prf:negation-disjunction-functionally-complete-propositional-logic
